%% bioETH-PRS: Confidential Polygenic Risk Scoring with Smart Contracts
%% Academic paper — privacy, genomics, blockchain

\documentclass[10pt]{article}

\usepackage[T1]{fontenc}
\usepackage[utf8]{inputenc}
\usepackage[a4paper, top=2.0cm, bottom=2.4cm, left=1.5cm, right=1.5cm, columnsep=0.65cm]{geometry}
\usepackage{amsmath}
\usepackage{amssymb}
\usepackage{amsthm}
\usepackage{booktabs}
\usepackage{hyperref}
\usepackage{microtype}
\usepackage{graphicx}
\usepackage{url}
\usepackage{xcolor}
\usepackage{array}
\usepackage{tabularx}
\usepackage{titlesec}
\usepackage{caption}
\usepackage{enumitem}
\usepackage{natbib}
\usepackage{fancyhdr}

\hypersetup{colorlinks=true,linkcolor=blue!60!black,citecolor=blue!60!black,urlcolor=blue!60!black}

\titleformat{\section}{\normalfont\large\bfseries\scshape}{\thesection}{1em}{}
\titleformat{\subsection}{\normalfont\normalsize\bfseries}{\thesubsection}{1em}{}
\titleformat{\subsubsection}{\normalfont\normalsize\itbfseries}{\thesubsubsection}{1em}{}

\setlist[itemize]{leftmargin=*, topsep=2pt, itemsep=1pt}
\setlist[enumerate]{leftmargin=*, topsep=2pt, itemsep=1pt}

\captionsetup{font=small, labelfont=bf}
\raggedbottom
\setlength{\emergencystretch}{2em}
\setlength{\parindent}{0pt}
\setlength{\parskip}{4pt plus 1pt minus 1pt}
\hbadness=10000
\vbadness=10000
\hfuzz=4pt
\vfuzz=4pt

\newcolumntype{Y}{>{\raggedright\arraybackslash}X}

\newtheorem{proposition}{Proposition}

%% Footer setup
\pagestyle{fancy}
\fancyhf{}
\fancyfoot[R]{\small\thepage}
\renewcommand{\headrulewidth}{0pt}
\renewcommand{\footrulewidth}{0.4pt}
%% (no footoffset needed)

%% ============================================================
\title{\textbf{bioETH-PRS: Confidential Polygenic Risk Scoring with\\[2pt]
       Smart Contracts on an FHE-Enabled Blockchain}}

\author{%
  \begin{minipage}{0.96\textwidth}
  \raggedright
  \textbf{Kimon Antonios Provatas$^{1,+,*}$, Christos Galanopoulos$^{1,+}$, Ilias Georgakopoulos-Soares$^{1,*}$}\\[2pt]
  \normalsize $^{1}$Division of Pharmacology and Toxicology, College of Pharmacy, The University of Texas at Austin, Dell Pediatric Research Institute, Austin, TX, USA\\[2pt]
  \small $^{+}$Co-first authors \quad $^{*}$Corresponding authors: \texttt{kap4722@my.utexas.edu}; \texttt{ilias@austin.utexas.edu}\\[2pt]
  \small\textit{Blockchain $\cdot$ Genomic Privacy $\cdot$ FHE $\cdot$ Smart Contracts $\cdot$ Risk Categories}
  \end{minipage}
}
\date{}

\begin{document}
\maketitle
\thispagestyle{fancy}
\vspace{-0.8em}

%% ============================================================
%% GRAPHICAL ABSTRACT (full-width, before abstract)
\begin{center}
  \centering
  \includegraphics[width=0.70\textwidth,keepaspectratio]{figures/graphical_abstract.png}
  \captionof{figure}{\textbf{Graphical Abstract.} bioETH-PRS uses smart contracts to coordinate
  an encrypted PRS calculation instead of relying on a designated evaluator. The system
  still depends on the blockchain and on fhEVM services that perform encrypted
  calculations and decrypt results for the intended requester. We report a public-weight 100-SNP calculation
  on Sepolia, exact agreement for 200 calculations in a local simulation, and
  35.4--37.2\% lower gas for the Streaming method in the local simulation.}
  \label{fig:graphical_abstract}
\end{center}
\vspace{-0.8em}

%% ============================================================
\begin{abstract}
\noindent
Polygenic risk scores (PRSs) aggregate genetic effect estimates to predict disease
susceptibility.  Calculating a PRS with a conventional external service can require
sharing raw genotype data.  HEPRS, the homomorphic-encryption system used for comparison
in this study, requires a designated evaluator.  We present \textbf{bioETH-PRS}, which uses publicly auditable
smart contracts to coordinate the calculation on a blockchain that supports Fully
Homomorphic Encryption (fhEVM).  This reduces reliance on a single evaluator, but the
system still depends on the smart contracts, blockchain, and fhEVM computation and
decryption services.
Before encryption, the input is checked for valid genotype values and aligned to the
model's effect allele.  bioETH-PRS converts signed weights to nonnegative integers
without overflow.  The model provider can require a randomized risk category instead
of the raw score.  One public-weight 100-SNP calculation completed on
Sepolia in 25 transactions and used 20.710271 million gas.  It took 269,320 ms from
submission to result and 8,081 ms to decrypt; the encoded score of 758,685 matched an
independent calculation exactly.  One private-weight 100-SNP calculation was evaluated
only in a local simulation.  Across 200 local public-weight calculations with
100--5,000 variants, every decoded
score matched the independent calculation because the study weights were represented
without rounding.  The Streaming method processes each group without permanently storing
its inputs; it used 35.4--37.2\% less local gas than the Classic method (stored inputs).
This study evaluates additive PRS models containing
up to 5,000 variants; we did not evaluate genome-wide or clinical use.
\end{abstract}

%% ============================================================
\twocolumn[{
\section*{Key Points}
\begin{itemize}
  \item Smart contracts reduce reliance on a single designated evaluator, but
        the system still depends on the blockchain and fhEVM services.
  \item A public-weight 100-SNP calculation completed on Sepolia.  In the local simulation,
        we evaluated public-weight calculations with 100--5,000 variants and one
        private-weight 100-SNP calculation.
  \item All 200 locally decoded public-weight scores matched Equation~\ref{eq:prs} exactly because
        the study weights were represented without rounding.
  \item The Streaming method used 35.4--37.2\% less gas than the Classic method (stored inputs) in the
        local simulation.
  \item Randomized risk categories and query limits make exact recovery of private weights harder,
        but they do not provide differential privacy or completely hide the weights.
\end{itemize}
}]

%% ============================================================
\section{Introduction}
\label{sec:intro}

Advances in genome sequencing have made polygenic risk scores (PRSs) a central tool
in precision medicine \citep{wray2021polygenic,inouye2018genomic}.  A PRS aggregates
the dosage-weighted effects of genetic variants:
\begin{equation}
  \mathrm{PRS} = \sum_{i=1}^{N} g_i \cdot \beta_i,
  \label{eq:prs}
\end{equation}
where $g_i \in \{0,1,2\}$ is the dosage of the model-specified effect allele for variant
$i$ and $\beta_i$ is the corresponding GWAS effect weight.  Published polygenic scores in the PGS Catalog range
from small panels to genome-wide scores with large variant sets \citep{lambert2019polygenic}.
Using a conventional external service to calculate a PRS can require sending genotype
data outside the hospital, which raises privacy concerns because genomic sequences can
be identifying \citep{gymrek2013,erlich2014}.

\textbf{Two kinds of sensitive information.}
Two kinds of sensitive information require simultaneous protection: (1)
\emph{patient genotypes} ($g_i$), identifying, heritable, and permanent; and
(2) \emph{GWAS model weights} ($\beta_i$), derived from restricted research cohorts.
More generally, model interfaces and outputs can leak sensitive information through
model-inversion attacks \citep{fredrikson2015modelinversion}.
In a conventional plaintext service, the evaluator must decrypt both genotypes and
weights before the calculation can proceed.

\textbf{Why FHE is useful.}
Fully Homomorphic Encryption (FHE) permits arbitrary arithmetic on encrypted data
without decryption.  Knight et al.\ \citep{knight2026heprs} demonstrate this with
HEPRS, an open-source system that computes a 110,000-SNP schizophrenia PRS under
the CKKS scheme \citep{cheon2017ckks} with near-zero accuracy loss ($r > 0.999$,
MSE~$< 2.3 \times 10^{-6}$).  HEPRS, however, still uses a central evaluator:
it relies on a three-party trust model (client, modeler, evaluator), and a colluding
evaluator and client can reveal GWAS model weights.  Trust in the evaluator is an
assumption, not a protocol guarantee.  Furthermore, HEPRS requires 3--4~GB of RAM
per individual even at 110,000 SNPs, and output is a raw score---there is no built-in
mechanism to prevent score accumulation or model probing.

\textbf{Our contribution.}
We propose \textbf{bioETH-PRS}, in which publicly auditable smart contracts coordinate
the encrypted calculation and control how results are released.  A designated
evaluator is no longer required.  However, the system still relies
on the contracts, blockchain consensus, and the fhEVM computation and decryption
services (Figure~\ref{fig:graphical_abstract}).  The blockchain records the contract
steps, but it does not by itself prove that the encrypted arithmetic was correct.

\smallskip\noindent\textbf{Distinctions from HEPRS.}
Although both systems calculate the same encrypted PRS, they use different methods.
First, bioETH-PRS converts signed decimal weights to unsigned integers for TFHE/fhEVM;
HEPRS uses CKKS approximate decimal arithmetic.  Second, four public contracts manage
who may use a sample, model publication, calculation, and result release.  Third, the
Classic method (stored inputs) keeps the encrypted SNPs for later calculation, whereas the Streaming
method processes each group immediately.  Fourth, a model can require a randomized
risk category and limit repeated queries instead of releasing the raw score.

\smallskip\noindent Our main contributions are:
\begin{itemize}
  \item Four smart contracts that separately manage who may use a sample, model
        publication, encrypted calculation, and result release.
  \item A three-step conversion that represents signed decimal weights as nonnegative
        integers, checks that the calculation stays within the allowed integer range,
        and selects a suitable scale before the model is published.
  \item The Classic method (stored inputs) and Streaming method for processing the encrypted inputs;
        the Streaming method used about 37\% less gas in the local simulation.
  \item A randomized risk category that can be required instead of releasing
        the raw score.
  \item Query limits and thresholds fixed by the model provider, evaluated with queries
        chosen after earlier results or in advance, several wallets and samples, and
        correlated SNP patterns.
  \item A public-weight 100-SNP calculation on Sepolia, public-weight calculations with 100--5,000
        variants in the local simulation, and one local private-weight 100-SNP calculation.
\end{itemize}

The intended use is research with additive PRS models containing up to 5,000 variants.  We did not
evaluate genome-wide or clinical use, production cost, or private weights on Sepolia.

%% ============================================================
\section{Background}
\label{sec:background}

\subsection{Polygenic Risk Scores}

A PRS is the weighted inner product of a patient's effect-allele dosage vector
$\mathbf{g} \in \{0,1,2\}^N$ and a GWAS effect weight vector $\boldsymbol{\beta}
\in \mathbb{R}^N$ (Eq.~\ref{eq:prs}).  Effect weights are typically estimated by
large-scale GWAS over tens of thousands of individuals and are signed floating-point
values in the range $[-0.5, +0.5]$.  Standard tools such as PLINK~2
\citep{chang2015plink} and LDpred2 \citep{prive2020ldpred2} compute PRSs in
milliseconds on plaintext data; FHE approaches trade runtime for privacy.

PRS models span a wide range in SNP inclusion depending on construction methodology.
Sparse models may include hundreds to a few thousand variants, while genome-wide
approaches incorporate tens of thousands to millions.  We evaluated bioETH-PRS with
models containing up to 5,000 variants.  These results do not establish genome-wide or
clinical use.

\subsection{Genotype preprocessing, QC, and model alignment}
\label{sec:preprocessing}

bioETH-PRS scores one individual against an already developed model.  Minor-allele
frequency and Hardy--Weinberg filtering are therefore cohort- and model-development
quality-control operations performed upstream when effect weights are derived.  The
scoring procedure instead checks missingness, genome build, variant identity and
order, allele orientation, and the dosage representation before any value is encrypted.

Only diploid hard-call dosages in $\{0,1,2\}$ are accepted.  A non-integer or
out-of-range dosage is rejected rather than clamped, because clamping silently changes
the score.  The model must state one of three rules for missing values: reject the
sample, use dosage 0, or use a supplied cohort mean rounded half away from zero to a
hard call in $\{0,1,2\}$.  An unstated zero is avoided because it would be
indistinguishable from a genuine homozygous-reference call.
The declared genotype build must equal the model build; a mismatch is fatal because the
build cannot be inferred from dosage values.  Variant identifiers and order are checked
element by element, not merely by vector length.  Duplicate identifiers, multiallelic
sites, and indels are rejected; this study evaluates biallelic SNP hard calls only.

Effect-allele alignment does not reveal encrypted weight magnitudes.  Publicly visible model
information supplies each variant identifier, genome build, effect allele, other allele,
and column order, and alignment runs locally before encryption.  For a biallelic,
non-palindromic SNP, a dosage already counting the effect allele is retained; a dosage
counting the other allele is transformed as $g_{\mathrm{effect}}=2-g$.  A compatible
strand complement is applied before the same decision.  Palindromic A/T or C/G pairs
without explicit strand resolution are rejected even on a literal allele match, because
the same label is consistent with opposite strand orientations.  Any remaining
incompatible pair is rejected.

The preparation record reports how many variants were matched, missing, imputed,
rejected, or changed for allele orientation.  The HEPRS data sets contain
numeric matrices without variant identifiers, builds, or allele labels and are therefore
treated as already aligned.  Their leading constant column has weight 0 and dosage 1,
so an $N$-variant data set contains $N+1$ encoded positions.  Examples with variant
metadata covered the build, order, missingness, and strand rules.

\subsection{Fully Homomorphic Encryption}

FHE schemes differ in the arithmetic they support and their performance.  Two are
directly relevant to this work.

\textbf{CKKS} \citep{cheon2017ckks} supports approximate arithmetic over real-valued data, is
naturally suited to decimal calculations, and underpins HEPRS.  Its available
calculation depth depends on the selected parameters and whether bootstrapping is used,
and approximation error accumulates during calculation.  The Lattigo library
\citep{lattigo} implements CKKS in Go with the parameter sets used
by HEPRS.  CKKS is computationally expensive: evaluating a 110,000-SNP PRS on a
single CPU node requires approximately 4.9~s per individual and up to 130~GB of RAM
for a cohort of 1,000 \citep{knight2026heprs}.

\textbf{TFHE} \citep{chillotti2020tfhe} supports bootstrapped operations on encrypted
values.  The fhEVM interface used here exposes exact unsigned-integer operations to
smart contracts.  Signed decimal weights must therefore be converted to unsigned
integers, motivating Section~\ref{sec:quantization}.

The security of TFHE is based on LWE/GLWE-family assumptions
\citep{chillotti2020tfhe}; CKKS uses related lattice assumptions.  These assumptions
are generally studied as candidates for security against quantum attacks, but the
present study does not evaluate long-term cryptographic parameter selection.

\subsection{Programmable Blockchain and fhEVM}

The Fully Homomorphic Ethereum Virtual Machine (fhEVM) extends the Ethereum
execution environment with TFHE precompile operations \citep{zamafhevm}.  Every
encrypted value is represented on-chain by a 32-byte opaque handle; actual
ciphertexts are managed by an off-chain coprocessor.  Smart contracts store only
handles and trigger FHE operations through coprocessor calls.

An on-chain contract records who may decrypt each result.  If the fhEVM services work
as specified, only the intended requester can decrypt it.  The blockchain records these decisions and calls, but it does not
independently verify the encrypted calculation or prevent a failure of the fhEVM
calculation and decryption services.

A per-transaction Homomorphic Computation Unit (HCU) budget limits FHE operations.
In the local simulation, the largest successful calculation chunk contained 21 SNPs
for both model types (Section~\ref{sec:evaluation}); the Sepolia limit remains
unmeasured and is not inferred from that result.  The HCU constraint necessitates
chunked computation strategies for large input vectors.

%% ============================================================
\section{System design}
\label{sec:design}

\subsection{System overview}

bioETH-PRS uses four smart contracts to move from sample registration to a risk result
(Figure~\ref{fig:arch}).  The registry records who may use a sample, the Model
contract stores the GWAS model and its result settings, the Calculation contract calculates
the encrypted score, and the Result contract can release a randomized risk category.  Separating
these roles makes the responsibility of each contract clear.

\begin{figure}[!htbp]
  \centering
  \includegraphics[width=\columnwidth,keepaspectratio]{figures/fig_architecture.png}
  \caption{System architecture. Four contracts record who may use a sample, store the model
  and its result settings, calculate the encrypted score, and return either the score or a
  randomized risk category. Public weights and model details are visible by design;
  genotype dosages and private weights remain encrypted during the calculation.
  Results are released only to the intended requester.}
  \label{fig:arch}
\end{figure}

\textbf{Sample registry.}
The registry stores references (IPFS or Arweave URIs) to encrypted genotype data held
off chain.  A requester registers a sample and chooses who may use it in a calculation.
The registry keeps this record but performs no encrypted arithmetic.

\textbf{Model contract.}
The model contract stores GWAS weights in groups.  Public weights are visible
integers; private weights are encrypted.  Public weights are encrypted before
multiplication in the present implementation.  The model record also stores the scale
factor, weight shift, score shift, and a cryptographic fingerprint of the input-preparation
record.  Models cannot be changed after publication.  Public-weight calculations with
100--5,000 variants and one private-weight 100-SNP calculation were evaluated locally.
Only the public-weight 100-SNP calculation was performed on Sepolia.

\textbf{Calculation contract.}
The calculation contract receives the encrypted SNPs in groups, calculates the weighted
sum, corrects the integer shifts described in Section~\ref{sec:quantization}, and
releases one result.  The final encrypted score can be decrypted only by the intended
requester.  Section~\ref{sec:execution} describes the two calculation methods.

\textbf{Result contract.}
The result contract adds a random value to the encrypted score before assigning Low,
Medium, or High risk.
The random value comes from the fhEVM and is not known before the transaction is mined.
The model provider can require this category result instead of allowing the requester to
decrypt the raw score.

\subsection{Comparison with HEPRS}

Table~\ref{tab:comparison} compares bioETH-PRS and HEPRS
\citep{knight2026heprs} by how they handle encrypted data, model size, runtime, memory,
deployment needs, and the result returned.  Both compute Equation~\ref{eq:prs}.  HEPRS reports encrypted CKKS
results for a much larger model with a designated evaluator.  bioETH-PRS reports a
smaller calculation coordinated by public smart contracts.  These contracts reduce
reliance on one evaluator, but the system still depends on the blockchain and fhEVM
services.

\begin{table*}[t]
\centering
\caption{Comparison with HEPRS. HEPRS values come from the published study;
bioETH-PRS values state whether they were observed on Sepolia, observed in the local
simulation, or not measured.}
\label{tab:comparison}
\footnotesize
\setlength{\tabcolsep}{4pt}
\begin{tabularx}{\textwidth}{@{}p{3.1cm}XX@{}}
\toprule
\textbf{Dimension} & \textbf{HEPRS} \citep{knight2026heprs} & \textbf{bioETH-PRS} \\
\midrule
How encrypted data are handled & encrypted client--modeler--evaluator system & encrypted inputs; contracts determine who receives the result \\
Designated evaluator & yes; evaluator executes CKKS & no designated evaluator; contracts coordinate the calculation \\
Remaining dependencies & three-party non-collusion and HE software/hardware & contracts, blockchain, fhEVM computation, and decryption services \\
Arithmetic scheme & CKKS approximate signed-real arithmetic & TFHE unsigned integers after fixed-point encoding \\
Encrypted variants evaluated & 110,000 in the published study & public: 100 on Sepolia and 100--5,000 locally; private: 100 locally \\
Runtime & approximately 4.9 s/person at 110,000 & 269.320 s to result plus 8.081 s to decrypt at 100 on Sepolia; local times are not comparable \\
Memory & 3--4 GB/person; up to 130 GB for 1,000 & not measured \\
Deployment needs & evaluator compute host and client/modeler setup & fhEVM blockchain, computation service, relayer, and decryption service \\
Result returned & raw score delivered to client; external controls optional & raw score or randomized risk category, as set by the model provider \\
Public information & model information managed off chain & model identifiers, alleles, order, result settings, and transactions are public; weights may be encrypted \\
\bottomrule
\end{tabularx}
\end{table*}

%% ============================================================
\section{Representing decimal weights as integers}
\label{sec:quantization}

\subsection{Why conversion is needed}

This implementation represents weights and scores as unsigned 64-bit integers.
GWAS weights $\beta_i$ are signed floating-point values;
a naive fixed-point encoding that multiplies by a scale factor $s$ and rounds
produces quantized weights $q_i = \mathrm{round}(s \cdot \beta_i) \in \mathbb{Z}$,
which may be negative.  Negative integers cannot be stored in an unsigned 64-bit
integer and would silently wrap under modular overflow, corrupting results without any
observable error.  A correct encoding must guarantee that all intermediate encrypted
values remain in the non-negative range $[0, 2^{64} - 1]$ for any genomically valid
input $g_i \in \{0, 1, 2\}$.

\subsection{Three-step conversion}

We convert the signed decimal weights to nonnegative integers in three steps
(Figure~\ref{fig:quantization}).  The integer shifts are reversible.  Conversion of the
original decimal weights is exact only when the selected scale represents those weights
without rounding, as it does for the data sets used here.

\begin{figure}[t]
  \centering
  \includegraphics[width=\columnwidth,keepaspectratio]{figures/fig_quantization.png}
  \caption{Original three-stage overview of the fixed-point conversion. The worked
  example is converted back to the additive PRS. Equations~\ref{eq:wshift}--\ref{eq:encode}
  give the exact half-away-from-zero rule, unsigned clamp, and intermediate-value order used
  in the implementation.}
  \label{fig:quantization}
\end{figure}

\textbf{Step 1: Scale.}
\begin{equation}
q_i = \operatorname{round}_{\mathrm{away}}(s \cdot \beta_i), \quad q_i \in \mathbb{Z},
\end{equation}
where $s$ is a model-specific scale factor selected by the quantization procedure
(Section~\ref{sec:advisor}) and exact half ties are rounded away from zero.
The study data sets require no tie breaking: all weights have at most six decimal
places and the selected scales are integer multiples of $10^6$, so their quantization
is exact.

\textbf{Step 2: Weight zero-point shift.}  Let
\begin{equation}
z_w = \max(0,-\min_i q_i), \quad u_i = q_i + z_w \;\geq\; 0 \;\;\forall i.
\label{eq:wshift}
\end{equation}
The clamp is required by the unsigned on-chain representation when all weights are
positive; algebraically, $z_w=0$ already preserves $u_i\geq0$.  The shifted weights
$u_i \geq 0$ are stored on-chain.  The raw accumulation
$\sum_i g_i u_i$ overestimates the true PRS by the constant $z_w \cdot \sum_i g_i$,
requiring correction.  The engine tracks the genotype sum
$G = \sum_i g_i$ as a second encrypted accumulator to enable this correction.

\textbf{Step 3: Score zero-point shift.}
Even after the weight correction, the corrected sum
$P = \sum_i g_i u_i - z_w G$
can be negative for patients with many risk-decreasing alleles.  We define:
\begin{align}
z_s &= -\sum_{i:\beta_i<0} 2 q_i \;=\; \sum_{i:\beta_i<0} 2|q_i|,\\
e   &= P + z_s \;\geq\; 0.
\end{align}
$e$ is the \emph{encoded score} stored encrypted on-chain.  The on-chain
computation is structured to avoid any intermediate negative value:
\begin{equation}
e = \bigl(\mathrm{partialSum} + z_s\bigr) - \bigl(z_w \cdot G\bigr),
\label{eq:encode}
\end{equation}
where $z_s$ is added before subtracting $z_w G$, ensuring the intermediate
value $(\mathrm{partialSum} + z_s)$ is always non-negative.

\textbf{Decoding.}  After decryption by the intended requester:
\begin{equation}
\mathrm{PRS} = \frac{e - z_s}{s}.
\end{equation}

\subsection{Worked Example}

Consider three genetic variants with scale $s=100$.  Their weights and effect-allele
dosages are
\[
\boldsymbol{\beta}=[-0.30,0.10,0.25],\qquad \mathbf{g}=[0,2,1].
\]
The plaintext target is
\[
\mathrm{PRS}=0(-0.30)+2(0.10)+1(0.25)=0.45.
\]

\begin{enumerate}
  \item Quantize: $\mathbf{q} = [-30, 10, 25]$.
  \item Weight shift: $z_w = 30$, $\mathbf{u} = [0, 40, 55]$.
  \item Accumulate: partialSum $= 0{\cdot}0 + 2{\cdot}40 + 1{\cdot}55 = 135$;
        $G = 0 + 2 + 1 = 3$.
  \item Correction: $P = 135 - 30{\cdot}3 = 45$.
  \item Score shift: $z_s = 2{\cdot}30 = 60$; $e = 45 + 60 = 105$.
  \item Decode: PRS $= (105 - 60) / 100 = 0.45$.
\end{enumerate}

At the system level, the data preparer first checks variant identifiers, genome build,
dosage validity, missing values, and effect-allele orientation. Only then are dosage values
and, for private models, weight magnitudes encrypted. The requester selects an authorized
sample and model, submits the encrypted groups for calculation, and receives either the
model-configured raw score or randomized category. The contracts record permissions and the
result recipient, but they do not establish that the encrypted values came from the stated
biological sample.

\subsection{Checking the integer range}
\label{sec:overflow}

\begin{proposition}
Let $N$ be an arbitrary SNP count, let $s$ be the fixed-point scale factor, and
let $M$ satisfy $|q_i| \leq sM$ for every quantized weight $q_i$.  Then a
sufficient worst-case condition for unsigned 64-bit safety of the full encoding
calculation is
\[
4 s M N \leq 2^{64} - 1.
\]
Equivalently,
\[
N \leq \left\lfloor \frac{2^{64} - 1}{4 s M} \right\rfloor.
\]
Under this condition, both the final encoded score $e$ and the intermediate
quantity $\mathrm{withOffset} = \mathrm{partialSum} + z_s$ in
Eq.~\ref{eq:encode} fit in $\texttt{uint64}$.
\end{proposition}

{\renewcommand{\qedsymbol}{}
\begin{proof}
Partition the quantized weights into positive, negative, and zero sets with
cardinalities $N_{+}$, $N_{-}$, and $N_{0}$, so that $N_{+}+N_{-}+N_{0}=N$.
Because $g_i \in \{0,1,2\}$ and $z_w=\max(0,-\min_i q_i) \leq sM$, the shifted weight
$u_i = q_i + z_w$ obeys
\[
u_i \leq
\begin{cases}
2 s M, & q_i > 0,\\
s M, & q_i \leq 0.
\end{cases}
\]
Hence the encrypted accumulation before correction satisfies
\[
\mathrm{partialSum} = \sum_i g_i u_i
\leq 4 s M N_{+} + 2 s M (N_{-}+N_{0}).
\]
The score zero-point satisfies
\[
z_s = -\sum_{i:q_i<0} 2 q_i \leq 2 s M N_{-}.
\]
Therefore the intermediate value in Eq.~\ref{eq:encode} is at most
\begin{align*}
\mathrm{withOffset} &= \mathrm{partialSum} + z_s \\
&\leq 4sMN_{+}+4sMN_{-}+2sMN_{0} \\
&\leq 4sMN.
\end{align*}
This is the worst-case upper bound for the intermediate sum.

For the final encoded score, observe that the weight-shift correction exactly
cancels the zero-point contribution:
\[
P = \mathrm{partialSum} - z_w G = \sum_i g_i q_i.
\]
Thus
\[
-2 s M N_{-} \leq P \leq 2 s M N_{+},
\]
and after adding $z_s$ we obtain
\[
0 \leq e = P + z_s \leq 2 s M (N_{+}+N_{-}) \leq 2 s M N.
\]
So the encoded score range itself fits in $\texttt{uint64}$ under the weaker bound
$2 s M N \leq 2^{64}-1$, but the actual implementation must also accommodate the
larger intermediate $\mathrm{withOffset}$.  A sufficient safety condition for all
intermediate and final values is
\[
4 s M N \leq 2^{64}-1.
\]
This condition was satisfied by every model evaluated in this study.  It is an
arithmetic safety check, not evidence that the system can process models larger than
the evaluated maximum of 5,000 variants.
\end{proof}}

\subsection{Selecting the scale factor}
\label{sec:advisor}

Before publishing a model, the quantization procedure evaluates the candidate scales
\[
s \in \{10^2,10^4,10^6,10^8,10^{10}\}
\]
against the actual GWAS weight distribution.  For each $s$, it computes the mean absolute error
(MAE) between the quantized dot product and the floating-point PRS for all
individuals in the study data set.  It compares three ranges:

\begin{itemize}
  \item \textbf{Low scale} ($s \approx 10^2$): 1--15\% MAE in these data.
  \item \textbf{Selected scale} ($s \approx 10^6$): zero error on these six-decimal-place
        data sets.  A candidate recommendation for similarly represented models, not a
        universal production default.
  \item \textbf{Higher scale} ($s \approx 10^8$): no improvement over the selected scale
        on the study data; the limiting factor is source data precision.
\end{itemize}

This calculation takes approximately 200~ms and should be performed before every model
publication.  The scale does not change gas use; the number of transactions needed to
submit the SNPs is the limiting factor.

%% ============================================================
\section{Calculation methods}
\label{sec:execution}

bioETH-PRS provides two methods that differ in whether encrypted SNPs are stored between
transactions (Figure~\ref{fig:protocol}).  Both methods produce the same score.

\begin{figure}[t]
  \centering
  \includegraphics[width=\columnwidth,keepaspectratio]{figures/fig_protocol.png}
  \caption{Calculation method comparison. The Classic method (stored inputs) uploads and stores SNPs before the
  score is calculated, so different parties can perform these steps. The Streaming
  method combines upload and calculation for one requester. Across the 100--5,000
  variants evaluated in the local simulation, streaming used 35.4--37.2\% less gas.}
  \label{fig:protocol}
\end{figure}

\subsection{Classic method (stored inputs)}

The requester uploads encrypted SNPs in groups of at most 32.  After upload, the score
is calculated in groups of 20 SNPs, just below the largest successful local group of
21.  The encrypted SNP references remain stored between transactions.  This allows a
different account, such as a relayer, to continue the calculation.  Under the encryption
assumptions the relayer cannot read the SNP values, but it can delay requests.  Its
transactions also reveal timing and addresses that are already visible on the
blockchain.  It therefore remains a dependency for availability and privacy.

\subsection{Streaming method}

The requester provides and processes each group of up to 20 encrypted SNPs in the same
transaction.  The values are multiplied by the model weights, added to the running sum,
and then discarded rather than stored for later transactions.

\textbf{Gas savings.}
The Classic method permanently stores an encrypted SNP reference and a record stating
who may use it.  The Streaming method avoids both records, which explains the observed
35.4--37.2\% reduction in local gas.  The remaining gas in this implementation comes
mainly from preparing each encrypted input and multiplying and adding encrypted values.

\textbf{Trade-off.}
The Classic method allows upload and calculation to be performed by different parties.
The Streaming method requires one requester to submit and process the inputs.

%% ============================================================
\section{Security assumptions and limits}
\label{sec:security}

\subsection{Attacker considered}

We assume an attacker can read public blockchain data, submit transactions, and make
repeated classification requests with chosen encrypted inputs
(Figure~\ref{fig:security}).  A requester who is allowed to use private weights may
submit values that do not come from the registered sample.  We assume that the attacker
cannot break TFHE encryption or bypass the contract rules about who may use data and
receive results.

\begin{figure}[t]
  \centering
  \includegraphics[width=\columnwidth,keepaspectratio]{figures/fig_security.png}
  \caption{Security assumptions and remaining limits. The contracts control who may
  request a calculation, how the result is released, and how often a sample may be
  queried. They do not prove that the encrypted SNPs came from the registered sample,
  that the model is clinically valid, or that repeated queries reveal no information
  about private weights.}
  \label{fig:security}
\end{figure}

\subsection{Source of the encrypted SNPs}

The registry checks whether a person may request a calculation, but it cannot check
whether the encrypted SNP values came from the registered biological sample.  The
contracts calculate a result from the values they receive, including values chosen by a
requester.  Restricting an attack to correlated SNP patterns is therefore an analysis
choice, not a restriction enforced by the system.

This study assumes that the patient, laboratory, or data holder prepares the genotype
data correctly before encryption, following Section~\ref{sec:preprocessing}.  The
stored sample record describes the genome build, variant order, preparation rules, and
a cryptographic fingerprint of that record.  It helps document how the input was prepared, but it does not prove
that the encrypted values came from the stated biological sample.

\begin{table*}[t]
\centering
\caption{Parts of the system on which privacy, correctness, availability, and the
sample or model record depend. A mark shows what could be affected if that part fails.}
\label{tab:trust_boundary}
\small
\begin{tabularx}{\textwidth}{@{}XccccX@{}}
\toprule
\textbf{Component} & \textbf{Privacy} & \textbf{Correctness} & \textbf{Availability} & \textbf{Record} & \textbf{Responsibility} \\
\midrule
Genotype provider/preprocessor &  & X &  & X & prepares the claimed sample and applies build, order, QC, and orientation rules \\
Model provider &  & X &  & X & publishes the intended weights, metadata, thresholds, and scientifically valid model \\
Smart contracts & X & X & X & X & apply the stated calculation and the rules about who may use data or receive results \\
Blockchain consensus &  & X & X & X & orders and finalizes the publicly auditable contract record \\
fhEVM coprocessor & X & X & X &  & performs TFHE operations correctly without exposing plaintext \\
Relayer & X &  & X &  & carries requests without revealing a released result to someone else \\
Decryption services & X & X & X &  & decrypt results only for the intended requester \\
\bottomrule
\end{tabularx}
\end{table*}

\subsection{Rules for protecting results}

Under the assumptions in Table~\ref{tab:trust_boundary}, the contracts apply five
rules to protect the results:
\begin{itemize}
  \item Contracts store opaque references to encrypted values, not plaintext SNPs.
        The requester creates the encrypted inputs, and the fhEVM services process them.
  \item Each encrypted output records who may decrypt it, and the contract gives this
        record only after the required calculation steps.
  \item The encoded PRS $e$ (Eq.~\ref{eq:encode}) is never publicly decryptable.
  \item A model may require the randomized Low, Medium, or High category instead of
        releasing the raw score to the requester.
  \item Each calculation releases one result; another result requires another query.
\end{itemize}
These rules do not establish where the submitted SNPs came from, whether the model is
clinically valid, or whether the external fhEVM services behave correctly.

\subsection{Randomized risk category}

Before assigning Low, Medium, or High risk, the Result contract adds a random
nonnegative integer smaller than $B$ to the encoded score:
\begin{equation}
\begin{aligned}
e_{\mathrm{randomized}} &= e + \nu,\\
\nu &\sim \operatorname{DiscreteUniform}\{0,1,\ldots,B-1\}.
\end{aligned}
\end{equation}
The fhEVM supplies $\nu$; the requester does not choose it.  The model provider sets $B$
and the two category thresholds before any request is made, and requesters cannot
change the thresholds between queries.  This makes it harder to narrow down an
encrypted score by repeatedly changing a comparison threshold.

All comparisons against classification thresholds operate on the randomized encrypted
score using encrypted boolean operations, avoiding any branching on hidden values:
\begin{align}
\hat{L} &= \mathbb{1}[e_{\mathrm{randomized}} < \tau_L], \notag\\
\hat{H} &= \mathbb{1}[e_{\mathrm{randomized}} \geq \tau_H], \notag\\
\mathrm{cat} &= \mathrm{select}\bigl(\hat{L},\;L,\;\mathrm{select}(\neg\hat{L}\wedge\neg\hat{H},\;M,\;H)\bigr).
\label{eq:classify}
\end{align}
The comparisons are carried out on encrypted values.  The two thresholds must be at
least $B$ apart.

\textbf{Interpretation and limits.}
This method does not provide an $(\varepsilon,\delta)$-differential privacy guarantee.
It uses one-sided noise, does not relate $B$ to a formal measure of how much one input
can change the score, and does not account for information from repeated queries.  A
formal privacy method would need to address each of these points.

\textbf{Bias correction.}
Because the random value is always nonnegative, its exact mean is $(B-1)/2$.  The
contract provides the integer threshold correction $B/2$; for $B=128$, this is 64,
while the exact mean is 63.5.  Scores close to a threshold remain uncertain.  In the
100-SNP study, 48 individuals outside this uncertainty range kept the same category.
The other two were each 64 units below a threshold, so their category could change.

\subsection{Analysis of repeated queries}

We examined whether a requester allowed to use a model with 20 private weights could
infer the weights by submitting chosen encrypted SNP values and observing the released
results.  The analysis included thresholds chosen by the requester, with each query
chosen after seeing earlier results or all queries chosen in advance; thresholds fixed
by the model provider; several wallets and samples; and correlated SNP patterns.  These
calculations were performed in the local simulation.  To make the analysis repeatable,
we used a fixed sequence of random integers from 0 through 127. A different
random sequence may give different exact recovery counts. The contract also limits the
number of requests for each model, requester, and sample over a set number of blocks.

\begin{table*}[t]
\centering
\caption{Results of the repeated-query analysis in the local simulation. ``Within
$B$'' is the number of weights estimated to within the randomization range.}
\label{tab:probing}
\small
\begin{tabularx}{\textwidth}{@{}Xrrrr@{}}
\toprule
\textbf{Information released and query choice} & \textbf{Queries} & \textbf{Pearson $r$} & \textbf{Sign acc.} & \textbf{Within $B$} \\
\midrule
Raw score available & 20 & 1.0000 & 100\% & 20/20 \\
Requester changes threshold after each result & 200 & 0.9999 & 100\% & 19/20 \\
Requester changes threshold after each result, first 20/20 & 260 & 1.0000 & 100\% & 20/20 \\
Requester-selected thresholds, queries chosen in advance & 320 & 0.6689 & 65\% & 0/20 \\
Thresholds fixed by model provider & 320 & 0.9388 & 70\% & 0/20 \\
Fixed thresholds with correlated SNP blocks & 320 & 0.0223 & 65\% & 0/20 \\
\bottomrule
\end{tabularx}
\end{table*}

In this analysis, changing the threshold after each result recovered 19 of 20 weights
within $B$ after 200 queries and first recovered all 20 after 260 queries. When all
queries were chosen in advance, none was recovered within $B$ after 320 queries
($r=0.6689$; 65\% of signs correct). With thresholds fixed by the model provider, none
of the 20 estimates was within $B$ after 320 queries.  However, Pearson $r=0.9388$ and
70\% sign accuracy show that the results still revealed some information about the
relative weights.  The fixed thresholds therefore made precise recovery harder but did
not completely hide the private weights.

Using several wallets did not increase the number of queries allowed for the same
sample, although different registered samples had separate limits.  For private
weights, the model provider decides who may use them.  When inputs were restricted to
correlated SNP blocks, the correlation between estimated and true weights fell to
$0.0223$.  This is not a reliable safeguard because requesters can still submit other
encrypted SNP values. Applying the studied limit of three calculations per 1,000 blocks
to 260 total queries gives calculated times of 288.9 hours with 12-second blocks or
48.1 hours with 2-second blocks for one registered sample. These are calculated examples,
not measured network times. The chosen $B=128$ was 1.34\% of the largest integer weight;
other models may require a different value.

%% ============================================================
\section{Evaluation}
\label{sec:evaluation}

\subsection{Where calculations were evaluated}

We report results from Sepolia and from a local contract simulation.  Sepolia provides
observed fhEVM transaction results.  The local simulation evaluates contract behavior,
transaction counts, and gas used in a local contract environment; it does not measure
real encrypted-computation time, Sepolia capacity, or production fees.

We evaluated public-weight local calculations with 100, 500, 1,000, and 5,000 SNPs
\citep{knight2026heprs}, and one private-weight local calculation with 100 SNPs.  Each
public data set also includes one leading constant position.  The largest successful calculation
chunk in the local simulation contained 21 SNPs for both public and private weights, and
we used a chunk size of 20.  The corresponding Sepolia limit was not measured.

On Sepolia (chain ID 11155111), a public-weight 100-SNP calculation completed in 25
transactions after four contract deployments.  It used 20.710271 million gas, took
269,320 ms from submission to result and 8,081 ms to decrypt, and returned an encoded
score of 758,685.  This score exactly matched the independent calculation.  The
private-weight 100-SNP calculation was evaluated locally but was not run on Sepolia.

\begin{table*}[t]
\centering
\caption{Sepolia and local results for the same public-weight 100-SNP calculation using the
Classic method (stored inputs), upload groups of 32, calculation groups of 10, and 25 transactions.
Sepolia used 10.42\% more gas than this one local simulation; this percentage is not a
general conversion between local and network results.}
\label{tab:live_validation}
\small
\begin{tabular}{lrrrrr}
\toprule
\textbf{Setting} & \textbf{Tx} & \textbf{Gas} & \textbf{Submit--result} & \textbf{Decrypt} & \textbf{Encoded score} \\
\midrule
Sepolia & 25 & 20.710271 M & 269,320 ms & 8,081 ms & 758,685 \\
Local simulation & 25 & 18.755864 M & 147 ms & 125 ms & 758,685 \\
\bottomrule
\end{tabular}
\end{table*}

\subsection{Agreement with an independent calculation}

We compared all 200 public-weight results with an independent calculation of
Equation~\ref{eq:prs}: 50 individuals at each of the
four model sizes.  The mean,
root-mean-square, and maximum absolute errors were all zero; all 200 scores matched
exactly; and Pearson $r=1$.  This supports the calculation for the studied inputs.

The exact agreement is not a general precision result.  The four data sets contain
6,600 variant weights plus four leading constant weights.  All 6,604 values have at
most six decimal places, and the selected scale represented them without rounding.
Other weight sets may show small rounding differences.

\begin{table}[!htbp]
\centering
\caption{Independent Equation~\ref{eq:prs} calculation versus decoded public-weight
bioETH-PRS in the local simulation. The study weights were represented without rounding.}
\label{tab:quantization}
\small
\begin{tabular}{rrrll}
\toprule
\textbf{SNPs} & \textbf{Scale $s$} & \textbf{$n$} & \textbf{MAE/RMSE/max} & \textbf{Exact} \\
\midrule
100   & $3\times10^6$ & 50 & 0/0/0 & 50/50 \\
500   & $3\times10^6$ & 50 & 0/0/0 & 50/50 \\
1,000 & $1\times10^6$ & 50 & 0/0/0 & 50/50 \\
5,000 & $1\times10^6$ & 50 & 0/0/0 & 50/50 \\
\bottomrule
\end{tabular}
\end{table}

For the randomized risk category at 100 SNPs, all 48 individuals outside the
threshold uncertainty range kept the same category.  Two individuals were inside the
range and are reported separately; one changed category in this calculation.

\subsection{Variant scale}

Table~\ref{tab:scale_evidence} reports the number of transactions used for one public
model and one sample.  The 100-SNP Sepolia result used the Classic method (stored inputs).  The local
results from 100 to 5,000 variants used the Streaming method; 5,000 variants was the
largest bioETH-PRS model evaluated.

\begin{table}[ht]
\centering
\caption{Transactions for the public-weight calculation on Sepolia and in the local
simulation.}
\label{tab:scale_evidence}
\small
\begin{tabular}{lrlr}
\toprule
\textbf{Setting} & \textbf{Variants} & \textbf{Method} & \textbf{Tx} \\
\midrule
Sepolia & 100 & Classic & 25 \\
Local simulation & 100 & Streaming & 15 \\
Local simulation & 500 & Streaming & 47 \\
Local simulation & 1,000 & Streaming & 88 \\
Local simulation & 5,000 & Streaming & 413 \\
\bottomrule
\end{tabular}
\end{table}

\subsection{Gas use as variant count increases}

Figure~\ref{fig:gas} and Table~\ref{tab:gas} report gas used by the two calculation
methods for four public-weight calculations in the local simulation.  Gas grows approximately
linearly over the measured 100--5,000-variant range; this is not a claim about scaling
on Sepolia.  Across this range, the Streaming method used 35.4--37.2\% less gas because
it did not store each encrypted SNP reference permanently.

\begin{figure}[!htbp]
  \centering
  \includegraphics[width=\columnwidth, keepaspectratio]{figures/fig_gas_scaling.png}
  \caption{Original gas-scaling visualization for the two local calculation methods.
  The graphic retains the original plotted values; Table~\ref{tab:gas} reports the final
  post-review measurements and is the numerical source for the revised comparison.}
  \label{fig:gas}
\end{figure}

\begin{table}[!htbp]
\centering
\caption{Gas used in one local calculation, rounded to 0.001 million. The
classic and streaming totals include model publication and result release but exclude
sample registration. Both methods used calculation chunks of 20 SNPs.}
\label{tab:gas}
\small
\begin{tabular}{rrrr}
\toprule
\textbf{SNPs} & \textbf{Classic} & \textbf{Streaming} & \textbf{Reduction} \\
\midrule
100   & 17.863 M & 11.538 M & 35.4\% \\
500   & 83.985 M & 53.063 M & 36.8\% \\
1,000 & 166.875 M & 105.013 M & 37.1\% \\
5,000 & 829.422 M & 520.487 M & 37.2\% \\
\bottomrule
\end{tabular}
\end{table}

\subsection{Gas used for each SNP}

Table~\ref{tab:persnp} breaks down the approximate gas used for each SNP in the local
simulation.  In the present implementation, preparing the encrypted input and
multiplying and adding encrypted values account for approximately 95,000--104,000 gas
per SNP before permanent storage is considered.  Changes to the fhEVM services or input
preparation could change these values.

\begin{table}[ht]
\centering
\caption{Approximate gas used for each SNP in the local simulation}
\label{tab:persnp}
\small
\setlength{\tabcolsep}{4pt}
\begin{tabularx}{\columnwidth}{@{}Yrr@{}}
\toprule
\textbf{Operation} & \textbf{Classic} & \textbf{Streaming} \\
\midrule
Prepare encrypted input          & $\sim$50,000 & $\sim$50,000 \\
Multiply encrypted values        & $\sim$27,000 & $\sim$27,000 \\
Add encrypted values             & $\sim$27,000 & $\sim$27,000 \\
Store encrypted SNP (classic only) & $\sim$25,000 & --- \\
Store who may use it (classic only) & $\sim$25,000 & --- \\
Read stored data and other overhead & $\sim$12,000 & --- \\
\midrule
\textbf{Total/SNP}               & $\approx$\textbf{166,000} & $\approx$\textbf{104,000} \\
\bottomrule
\end{tabularx}
\end{table}

\subsection{Latency}

Local-simulation times for each public-weight calculation were 157, 780, 1,672, and
8,819 ms for 100, 500, 1,000, and 5,000 variants, respectively.  Each data set also
includes one leading constant position.  The measured range was 1.554--1.763 ms per
position.  The time per position increased by about 13\% between the smallest and
largest local calculations.  These times combine the local arithmetic and
transaction overhead; they measure neither real TFHE evaluation nor network latency.
The separate Sepolia 100-SNP times appear in Table~\ref{tab:live_validation}
and are not compared as a speed ratio.

\subsection{Transactions, gas, and fee examples}

Table~\ref{tab:cost} reports the observed transaction use.  The Sepolia deployment used
four transactions and 5.892559 million gas, costing 0.0062781714 Sepolia test ETH; the
public-weight calculation used 25 transactions and 20.710271
million gas, costing 0.0252747648 test ETH.  These are test-network expenditures, not
production prices.  In the local simulation, a public-weight calculation with 100 SNPs
using the Streaming method used 15 transactions and 11.690 million gas, whereas the
private-weight calculation with 100 SNPs used 17 transactions and 23.508 million gas
(2.01$\times$).  The private-weight calculation is the one relevant to attempts to
recover the weights.  Decryption happens outside Ethereum and requires no additional
Ethereum transaction.

\begin{table}[ht]
\centering
\caption{Measured transaction use. The local streaming totals include sample
registration and come from a separate calculation from Table~\ref{tab:gas}; encrypted
inputs also cause small gas variation between otherwise identical calculations.}
\label{tab:cost}
\footnotesize
\setlength{\tabcolsep}{2pt}
\begin{tabularx}{\columnwidth}{@{}lXrr@{}}
\toprule
\textbf{Setting} & \textbf{Quantity} & \textbf{Tx} & \textbf{Gas} \\
\midrule
Sepolia & four-contract deployment & 4 & 5.892559 M \\
Sepolia & public-weight 100-SNP, Classic method & 25 & 20.710271 M \\
Local simulation & public-weight 100-SNP, Streaming method & 15 & 11.690 M \\
Local simulation & private-weight 100-SNP, Streaming method & 17 & 23.508 M \\
Local simulation & one-time category settings & 1 & 0.077 M \\
\bottomrule
\end{tabularx}
\end{table}

Recording a digital fingerprint for the model file and another for its preparation information
added 40,568 local gas per model. This added no encrypted arithmetic and did not change the
largest successful local calculation group.  The
randomized category settings require one additional model transaction (77,314 local gas).
Releasing a raw score used 169,898 local gas, whereas releasing a randomized category
used 432,230 local gas. These
are alternative ways to release one result and are not added together.

For illustration only, we multiplied three local gas measurements by hypothetical gas
prices: deployment (5,892,613 gas), the public streaming calculation (11,690,021 gas),
and the private streaming calculation (23,507,880 gas).  These are calculated estimates,
not observed costs.  At 1 gwei, the respective values are 0.005892613, 0.011690021,
and 0.02350788 ETH; at 30 gwei, they are 0.17677839, 0.35070063, and 0.7052364 ETH.
We make no USD conversion or conclusion about affordability or clinical or commercial use.

\subsection{Calculation checks and responsibilities}

The independent comparison summarized in Table~\ref{tab:quantization} shows agreement
with Equation~\ref{eq:prs} for the studied inputs; it is not a guarantee for every
model or deployment.  Table~\ref{tab:correctness_boundary} explains who is responsible
for each part of the result.

\begin{table*}[t]
\centering
\caption{Responsibilities and remaining limits.}
\label{tab:correctness_boundary}
\small
\begin{tabularx}{\textwidth}{@{}p{3.2cm}XX@{}}
\toprule
\textbf{Actor/component} & \textbf{Responsibility in this study} & \textbf{Remaining limit} \\
\midrule
Data preparer & check genome build, variant order, dosage validity, missing values, and effect-allele alignment & the contracts do not prove that the encrypted SNPs came from the stated biological sample \\
Model provider & supply the weights, model details, thresholds, and scientific justification & calibration or performance in populations that were not studied \\
Smart contracts & apply the stated integer calculation and rules about who may use data or receive results & the contracts cannot verify input claims or external services \\
fhEVM services & perform the encrypted arithmetic and release the result to the intended requester & blockchain consensus does not independently verify the encrypted arithmetic \\
Independent calculation & compare Equation~\ref{eq:prs} with 200 studied inputs & agreement is not a proof for other models, inputs, or deployments \\
Requester & inspect the published model details and blockchain transactions & clinical interpretation beyond the published model \\
\bottomrule
\end{tabularx}
\end{table*}

The contracts check who may receive a result, allow only one result per calculation,
keep the thresholds fixed by the model provider, require the minimum threshold gap,
apply query limits, and stop a private-weight calculation when the model provider no
longer allows it.  They do not establish the
source of the biological sample, clinical validity, calibration, or performance in
other populations.

%% ============================================================
\section{Who may use the data and calculation order}
\label{sec:data_use}

\subsection{Encrypted values and released results}

Each encrypted value has a contract record stating who may use or decrypt it.  In the
Classic method (stored inputs), the SNP handles and these records are stored for later
transactions.  In the Streaming method, they are kept only for the current transaction
and then discarded.  Avoiding these two permanent records for each SNP explains the
lower gas use of the Streaming method.  The encoded score is released only to the
intended requester.  If the model requires a risk category, only the randomized Low,
Medium, or High result is made public.

\subsection{Order of the calculation}

The contracts require the operations to occur in order: create a calculation, provide
the encrypted SNPs, calculate the weighted sum, and release one result.  The Classic and
Streaming methods cannot be mixed within the same calculation.

\subsection{Who may use encrypted weights}

For private weights, the model provider chooses which requesters may use the model.  The
contract checks this choice whenever a calculation group is processed.  If the provider
no longer allows a requester to use the model, the next group is rejected.  The contract
also checks who may use the registered sample when the calculation is created.

%% ============================================================
\section{Discussion}
\label{sec:discussion}

We evaluated bioETH-PRS with additive PRS models containing up to 5,000 variants.  Publicly
auditable smart contracts reduce reliance on a single designated evaluator, but
the system still depends on the blockchain and on fhEVM services that perform encrypted
calculations, submit transactions, and decrypt results for the intended requester.  A public-weight 100-SNP
calculation completed on Sepolia.  Public-weight calculations with 100--5,000 variants
and one private-weight 100-SNP calculation were evaluated in the local simulation.
These results do not establish genome-wide or clinical use.

\subsection{Comparison with HEPRS}

bioETH-PRS and HEPRS \citep{knight2026heprs} address different questions and are best
understood as complementary systems.

HEPRS reports encrypted PRS calculation for a substantially larger model and measures
CKKS performance under a three-party non-collusion assumption.  CKKS represents
approximate signed decimal values without the integer conversion used here.  HEPRS
therefore provides the relevant results for large variant counts and measured FHE
performance.

bioETH-PRS instead studies whether public smart contracts can coordinate the calculation
and control result release without a designated evaluator.  Its model
settings and transaction history are public, while the system still relies on the fhEVM
services.  It also requires many blockchain transactions.  The two systems answer
different questions, and these results do not show that either is broadly superior.

\subsection{Limitations}

\textbf{Variant scale.}
The largest successful calculation chunk in the local simulation contained 21 SNPs for
both public and private weights; the Sepolia limit remains unknown.  A public-weight
5,000-variant calculation required 413 local transactions.  We did not evaluate larger
bioETH-PRS models or production feasibility.

\textbf{Source of the SNP data.}
As specified in Section~\ref{sec:security}, being allowed to request a calculation does
not prove that the encrypted SNP values came from the registered sample.

\textbf{URI observability.}
Sample data URIs are visible on the blockchain.  Any network participant can read them,
even if the application limits who may use the sample.  Storing only a cryptographic
fingerprint of the URI would reduce this exposure.

\textbf{One-sided randomization and bias.}
Adding an integer chosen uniformly from $0$ through $B-1$ increases the score by
$(B-1)/2$ on average and creates an uncertainty range near each threshold.  This
randomized category result is not
differential privacy because the random value is one-sided, is not formally calibrated
to score sensitivity, and is not analyzed across repeated queries.

\textbf{Production cost.}
We report gas used by the contracts, Sepolia test-ETH expenditure, and calculated fee
examples using hypothetical gas prices.  Production fee schedules, USD cost, memory
use, and operational throughput were not measured, so whether the system is affordable
or practical for clinical or commercial use remains unknown.

\subsection{Future work}

Future work should reduce the number of blockchain transactions and evaluate the
private-weight 100-SNP calculation on Sepolia.  Establishing the source of encrypted
SNPs will require a signed laboratory record or a privacy-preserving proof that links
the encrypted values to the registered sample.  A formal privacy guarantee will require
a precise definition of which data sets are compared, how much one input can change the
score, how noise is chosen, and what repeated queries reveal.  Later studies could also
evaluate LD-aware scores \citep{prive2020ldpred2}, decryption shared among several parties, model versioning,
and reuse of encrypted SNP values across models.

%% ============================================================
\section{Related Work}
\label{sec:related}

Privacy-preserving genomic computation has been pursued through multiple
cryptographic paradigms.  Kim and Lauter \citep{kim2015} applied homomorphic
encryption schemes to compute GWAS statistics for up to 5,000 sequences, but did not
extend to PRS computation.  Blatt et al.\ \citep{blatt2020} built a CKKS-based
privacy-preserving GWAS calculation for large-scale studies, noting that PRS
computation from decrypted summary statistics was future work.  McLaren et al.\
\citep{mclaren2016} demonstrated privacy-preserving genomic testing in an HIV
clinical-care model using homomorphic encryption, but did not address the full PRS
calculation.  Raisaro et al.\ \citep{raisaro2019} introduced MedCo for secure,
privacy-preserving exploration of distributed clinical and genomic data.

Knight et al.\ \citep{knight2026heprs} (HEPRS) is the closest prior work, providing
the first complete FHE calculation of PRS with real clinical data.  HEPRS achieves
$r > 0.999$ correlation with plaintext scores on a 110,000-SNP schizophrenia
model and demonstrates practical feasibility on commodity hardware.  bioETH-PRS
builds directly on the HEPRS computational framework and uses the same GWAS study
data, studying a different trade-off in which public contracts coordinate the
calculation while infrastructure trust remains.

On the blockchain side, prior work on privacy-preserving smart contracts has
focused on anonymous payments and related protocol work
\citep{sasson2014zerocash,pertsev2019tornado} and
zero-knowledge proof-based computation \citep{ben2018starkware}.  Zero-knowledge proofs
can verify a calculation while keeping designated witness data private, whereas FHE
allows calculations directly on encrypted data.  FHE on the blockchain, typified by the fhEVM
framework \citep{zamafhevm}, combines confidential operations with public smart-contract
records while depending on external encrypted-computation services.  bioETH-PRS applies
this framework to polygenic risk score calculation.

%% ============================================================
\section{Conclusion}
\label{sec:conclusion}

We presented bioETH-PRS, a protocol for confidential polygenic risk
scoring using TFHE-based fully homomorphic encryption on a programmable blockchain.
Publicly auditable smart contracts coordinate the calculation without a designated
evaluator.  This reduces reliance on one evaluator but does not
remove trust: the result still depends on genotype preparation, model validity, the
contracts, blockchain, and fhEVM services that perform encrypted calculations and
decrypt results for the intended requester.

One public-weight 100-SNP calculation completed on Sepolia and matched the independent
calculation exactly.  Public-weight calculations with 100--5,000 variants and one
private-weight 100-SNP calculation were evaluated only in the local simulation.  All
200 public-weight local scores matched
Equation~\ref{eq:prs} exactly because the study weights were represented without
rounding.  The Streaming method used 35.4--37.2\% less local gas.  In the repeated-query
analysis, thresholds fixed by the model provider and query limits made precise weight
recovery harder,
but some information about relative weights remained.  The randomized category result
does not provide differential privacy.

The results apply to additive PRS models containing up to 5,000 variants.  This study
does not establish genome-wide or clinical use, production affordability,
private-weight calculation on Sepolia, or the largest calculation step that Sepolia can
support.  HEPRS demonstrates encrypted calculation at a larger scale; bioETH-PRS studies
how smart contracts can coordinate a smaller calculation and control result release.

%% ============================================================
\section*{Biographical Note}
Kimon Antonios Provatas, Christos Galanopoulos, and Ilias Georgakopoulos-Soares are researchers at The University of Texas at Austin studying pharmacology, genomics, and privacy-preserving biomedical computation at the Dell Pediatric Research Institute.

%% ============================================================
\section*{Data Availability Statement}
No new datasets were generated or analyzed in this study.

%% ============================================================
\section*{Code Availability Statement}
All relevant code and results can be found on GitHub at \url{https://github.com/Georgakopoulos-Soares-lab/bioETH-PRS}.

%% ============================================================
\bibliographystyle{abbrvnat}
\bibliography{bioeth_prs}

\end{document}
